\subsection{The eighteen infrastructure files, one by one}
\label{sec:infra-table}

Table~\ref{tab:infra} is the concrete form of the claim made above. It lists every file of
the infrastructure layer with its size, what it builds, the closest library artefact, and a
verdict. The verdicts are: \emph{free} if the project's statement is the library's statement
and a one-line restatement suffices; \emph{partial} if the library covers part of the file or
covers it under hypotheses that have to be discharged; \emph{keep} if the object is
project-specific and a library would not want it; and \emph{notion} if the library has a
theory of the same name but a different object, so that replacement is a redesign rather
than a substitution.

\begin{center}\scriptsize
\begin{longtable}{@{}p{0.235\textwidth}rrp{0.205\textwidth}p{0.215\textwidth}p{0.09\textwidth}@{}}
\caption{The eighteen infrastructure files, what they build, the closest library artefact and
a verdict.}\label{tab:infra}\\
\toprule
file & lines & decls & what it builds & library counterpart & verdict \\
\midrule
\endfirsthead
\toprule
file & lines & decls & what it builds & library counterpart & verdict \\
\midrule
\endhead
\file{NavierStokes/}\newline\file{SmoothCutoffs} & 309 & 35 & smooth cutoffs on balls and
  strips, all-order derivative bounds, time switch & \Mathlib{} \lean{ContDiffBump},
  \lean{expNegInvGlue}; \TauCeti{} \file{BumpFunction/Cutoff} & free \\
\file{NavierStokes/}\newline\file{FlatCutoff} & 204 & 26 & the edge $\exp(-c/x^2)$ and its
  inverse-power family $x^{-\ell}\exp(-c/x^2)$ & \Mathlib{} \lean{expNegInvGlue}
  (\emph{without} the $c$-family) & partial \\
\file{NavierStokes/}\newline\file{FlatZeroExtension} & 243 & 21 & zero extension of a jet
  across such an edge & none & keep \\
\file{NavierStokes/}\newline\file{JetBounds} & 311 & 27 & finite-order bounds on
  \lean{iteratedFDeriv}, Leibniz profile & \Mathlib{} \lean{norm\_iteratedFDerivWithin\_mul\_le};
  \TauCeti{} \file{IteratedFDeriv/Basic} & partial \\
\file{NavierStokes/}\newline\file{WeightedClasses} & 517 & 37 & the weighted band-indexed
  coefficient classes $W_\alpha,M_\alpha,S_\alpha$ & none & keep \\
\file{NavierStokes/}\newline\file{ClosedIntervalCk} & 337 & 34 & the Banach space
  $C^k([-1,1];E)$ as a closed submodule of jet data & none (\TauCeti{}
  \file{IteratedDerivWithin} is pointwise, not a space) & keep \\
\file{NavierStokes/}\newline\file{ClosedIntervalJetAlgebra} & 293 & 31 & polynomial
  expressions in jets and their bounds & \Mathlib{} \lean{MvPolynomial} plus a bespoke
  interval-arithmetic interpreter & partial \\
\file{NavierStokes/}\newline\file{SmoothParameterIntegral} & 315 & 20 & all-orders
  differentiation under an integral over an arbitrary measure space & \TauCeti{}
  \file{Calculus/ParametricIntegral} covers the compact-parameter case & partial \\
\file{NavierStokes/}\newline\file{SchwartzCompactApproximation} & 58 & 4 & quantitative
  density of compactly supported smooth functions in $\mathcal S$, and the consequence that
  compact tests determine a tempered distribution & \TauCeti{} \lean{SchwartzMap.dense\_hasCompactSupport}
  & \textbf{done} \\
\file{Euler/}\newline\file{OrdinaryMollifier} & 164 & 26 & $L^2$-valued mollification along
  an isometric action, with rate & \TauCeti{} \file{Sobolev/Mollification} proves the
  weak-derivative identity instead & notion \\
\file{Euler/}\newline\file{CylinderSobolevSpace} & 145 & 25 & Sobolev spaces of all orders on
  a cylinder as a lattice of derivative words, complete by construction & \TauCeti{}
  \file{Sobolev/*} and \Mathlib{} \lean{BesselPotentialSpace} are Fourier-based and
  predicate-level & notion \\
\file{Euler/}\newline\file{SobolevWordLevel} & 54 & 5 & the level structure of that lattice &
  none & keep \\
\file{Euler/}\newline\file{SobolevSmoothApproximation} & 109 & 11 & smooth approximation in
  the word-lattice topology & \TauCeti{} \file{Sobolev/Mollification/Interior} (different
  carrier) & notion \\
\file{Euler/}\newline\file{SobolevHeatKernel} & 156 & 11 & heat gain of $k$ derivatives with
  an explicit parabolic majorant & \TauCeti{} \file{PDE/HeatKernel/Smoothing} (different
  carrier) & partial \\
\file{Euler/}\newline\file{CompactParameterIntegral} & 74 & 5 & $C^\infty$ parametric
  integral over a compact interval, from \emph{local} boundedness of the derivatives &
  \TauCeti{} \file{Calculus/ParametricIntegral} asks for bounds \emph{uniform in $x$}, which
  joint smoothness does not give; only the \lean{ContDiffOn} variant applies & partial \\
\file{Euler/}\newline\file{LpBochnerRealization} & 88 & 6 & $L^2(\mu;L^2(\nu))\cong
  L^2(\mu\otimes\nu)$ with equality of norms & \Mathlib{} \lean{MeasureTheory.Lp},
  \lean{integrable\_prod\_iff}; the identification is a short consequence & free \\
\file{Euler/}\newline\file{FiniteGradeAlgebra} & 109 & 10 & a truncated graded algebra with
  evaluation and a convolution inverse & \Mathlib{} \lean{Polynomial} evaluation and
  \lean{Finset.sum\_range} algebra & partial \\
\file{NavierStokes/}\newline\file{AppendixHeatResults} & 177 & 12 & rising factorials,
  $\Gamma(b+n)$ recursion, endpoint jets of the radial heat profile & \Mathlib{}
  \lean{ascPochhammer}, \lean{Real.Gamma\_add\_one}; \lean{rising} is now
  \lean{(ascPochhammer\ \mathbb{R}\ n).eval\ b} & \textbf{done} \\
\midrule
total & 3{,}674 & 387 & & & \\
\bottomrule
\end{longtable}
\end{center}
\noindent The verdict \lean{free} means a one-line restatement suffices; \lean{partial}
means the library covers part of the file; \lean{notion} means the library has a different
object of the same name; \lean{keep} means no library would want it.

\paragraph{What the table says.} Five substitutions have since been made, each verified by a green build.
\file{AppendixHeatResults} is done: its rising factorial \emph{is} Mathlib's
\lean{ascPochhammer} evaluated, with \lean{rising\_succ} and \lean{rising\_pos} now one-line
applications and only the monotonicity lemmas still by induction. In \file{R3/PairedKernelBound}
a duplicate of \lean{lpNorm\_one\_eq\_integral\_norm} was deleted, the file using the copy in
\file{R3/LpNormTools}. And in \file{R3/PressureFlux} the identity
$\|\mathrm{d}f(\lambda y.\,y{:}\mathbb{C})\| = \|\mathrm{d}f\|$ had been proved by an
antisymmetry argument over the operator norm in fifteen lines; it is now six, because
\Mathlib{} has \lean{LinearIsometry.norm\_toContinuousLinearMap\_comp} and
\lean{Complex.ofRealCLM} is \lean{Complex.ofRealLI.toContinuousLinearMap} by \lean{rfl}. The
net line change of those three is $+3$: what was removed is a bespoke proof, not a
statement.

The fifth is the one with volume. \file{R3/SchwartzCompactApproximation} reconstructed
density of compactly supported smooth functions in Schwartz space from a cutoff family and a
seminorm estimate---the whole file was 235 lines and 20 declarations---and nothing in the
library used the construction. \TauCeti{} proves the statement for the same object
(\lean{SchwartzMap.dense\_hasCompactSupport}, in \file{SchwartzSpace/Cutoff}), so the file
is now the import plus the two closed-set consequences of density: 235 $\to$ 58 lines,
20 $\to$ 4 declarations. Two further project theorems, \lean{image\_exp\_Iic} and
\lean{image\_exp\_Ioi}/\lean{Ioc}, were Mathlib's \lean{Real.image\_exp\_*} verbatim and are
gone, with their eleven uses repointed.

Finding these was mechanical rather than by reading. The duplicate probe of
Section~\ref{sec:dup} was run over the whole \file{NavierStokes} library instead of the
22-module sample, and it reports \emph{seven} theorem-duplicate groups involving a library
theorem: the two deliberate \TauCeti{} bridge aliases, one project--project pair, and the
four project theorems replaced above. That is the ceiling of statement-level substitution in
this library, and it is small because the library genuinely does not restate \Mathlib{}. The
Schwartz file is the one case where the duplication was of a \emph{construction} rather than
of a statement, and the probe cannot see those; it took reading the table in this section to
find it. That leaves two done, one free, eight partial, three notion-level and four to keep. (The parametric-integral file was recorded as
free in an earlier draft of this table; checking its hypotheses against \TauCeti's showed
that the library theorem needs bounds uniform in the differentiation variable, which the
project's joint-smoothness hypothesis does not supply, so it moved to \emph{partial}.) Nothing in the layer is large: the whole of it is
3{,}674 lines, which is 0.7\% of the development. The conclusion is that a library migration
of this repository is worth doing for the two free files and the eight partial ones, and is \emph{not} worth doing for the
word-lattice Sobolev spaces: those exist because the classical object was not available in
the library at the time, and replacing them now would change the meaning of the estimates
that the Euler development is made of. The honest intermediate step is the bridge lemma
\lean{SobolevSpace\ q\ →\ MemSobolev\ q\ 2}, which would connect the project's carrier to
\Mathlib's predicate and let the two theories be used side by side; it is not proved here.

\paragraph{Is any of this a special case of a general library lemma?} The duplicate probe
answers the question only for theorems whose \emph{statements} coincide. A project lemma can
also be a special case of a strictly more general library lemma, and that would be worth
replacing. To test it we asked the library directly: for each theorem, replace its proof by
\lean{exact?} \emph{in its own context} (no re-printing of types, so nothing is lost in
translation) and see whether a library term closes the goal.

The sample is the population where such a replacement is plausible: the \textbf{509}
\file{NavierStokes} theorems whose statement mentions \emph{no} project-defined constant, of
which the \textbf{173} in the 28 largest such files were run. The instrument was validated
first, on \lean{example : 0 < 2 * Real.pi := by exact?} and
\lean{example (x : ℝ) : x + 0 = x := by exact?}, which must --- and do --- report
\lean{Real.two\_pi\_pos} and \lean{AddMonoid.add\_zero}. (The first run of the harness looked
for the wrong line of \lean{exact?}'s two-line output and reported 0 of 173; the control is
what caught it.)

\textbf{9 of 173} are closed by a suggestion, and of those nine only \textbf{four} name a
library declaration:

\begin{center}\small
\begin{tabular}{@{}ll@{}}
\toprule
project theorem & library lemma that proves it \\
\midrule
\lean{real\_smul\_as\_complex\_mul} & \lean{Complex.real\_smul} \\
\lean{ae\_radial\_ne\_zero} & \lean{Measure.ae\_ne volume 0} \\
\lean{deriv\_smooth} & \lean{ContDiff.deriv'} \\
\lean{two\_le\_infty} & \lean{Std.LawfulOrderOrd.isGE\_compare} \\
\bottomrule
\end{tabular}
\end{center}

The other five suggestions are \emph{project} lemmas, which is the healthy case: the
development already factors them. All four replacements were applied; each saved one or two
lines, because each theorem was already a one-liner. The remaining 164 theorems are not
special cases of anything in the library: they are quantitative estimates about specific real
numbers, and \Mathlib{} has no general lemma to offer.

So the library overlap of this development is empty at the statement level (no theorem
duplicates a library theorem, checked against the whole of \Mathlib{}), and near-empty at the
level of generality (four one-line special cases out of 173 sampled). Whatever else accounts
for 341{,}000 lines, it is not library material.
